\begin{abstract}
Why do large-scale transformations routinely under-deliver, and why do 
``AI replacements'' appear plausible in domains once protected by long 
apprenticeships? Because we spent centuries training humans to be 
biological AI systems—and now the silicon versions have arrived to 
collect their inheritance.

This paper argues that both phenomena reflect the same structural 
constraint: an extraction--investment asymmetry. Institutions 
increasingly treat cognitive capability as a stock to be harvested 
(standardized, measured, and modularized) rather than a high-energy 
state that must be continuously maintained. The result is predictable 
brittleness under complexity and progressive deskilling toward tasks 
that are legible to automation.

We develop a thermodynamic framing as an \textit{analytic constraint} 
(not metaphysical equivalence): cognitive sovereignty depends on 
sustained energy allocation to synthesis and on resistance to 
extraction. We operationalize this with a sphere--vector geometric 
model, distinguishing integrative, context-dependent ``sphere'' 
capabilities from decomposable, assessable ``vector'' competencies. 
We trace how education and management systems privileged 
discretization, vectorization, and optimization—credits, competency 
frameworks, auditability—from historical formation regimes to 
contemporary micro-credentialing and Bologna-era standardization. We 
then examine contemporary deployment evidence (initiative abandonment, 
pilot-to-production brittleness, validated performance gaps, liability 
rulings, and governance failures) as mechanism-probes rather than 
prevalence claims.

The paper contributes (1) a unifying constraint linking educational 
degradation, organizational brittleness, and automation 
substitutability, (2) ``confession literature'' as a methodological 
resource for identifying standardization and extraction techniques, 
and (3) a research program with falsifiable probes and measurable 
indicators. The implication is not that AI ``cannot work,'' but that 
strategies optimizing for extractive legibility while underinvesting 
in maintenance, validation, and human capability preservation 
predictably degrade both people and systems. The window for capability 
preservation narrows as extraction accelerates. Constraints are 
non-negotiable; strategies are.
\end{abstract}